\chapter{Estabilidad y dinámica del conjunto rotor--cuerpo--drogue}\label{cap:estabilidad}
\usepgfplotslibrary{groupplots}

El conjunto es estable en todos los modos linealizados que se estudian. Hay tres puntos
débiles. Primero, el punto de diseño queda en el borde superior de la región de anillo de
vórtices. Segundo, el cuerpo gira por el rozamiento de los rodamientos: sin aletas llega
a cientos de rpm en medio minuto. Tercero, la pala de la configuración A tiene un margen
de divergencia torsional insuficiente a las \num{1505}~rpm del BEMT. Este capítulo usa
los valores refinados de los capítulos~\ref{cap:bemt} y~\ref{cap:estructura}:
$\Omega=\SI{157.6}{rad/s}$, $V_2=\SI{6.94}{m/s}$ con drogue y \SI{7.91}{m/s} sin él,
$T=\SI{3.58}{N}$, $C_R=0{,}97$, $\theta=\ang{-3}$, $\beta_0=\ang{6.1}$, número de Lock
$\gamma=11{,}5$, $I_b=\SI{4.35e-5}{kg.m^2}$, tope elástico a \ang{+15} y rodamientos
628/8-2Z (688ZZ). Los cálculos están en \texttt{analysis/estabilidad\_modos.py}.

\begin{resultado}[Resultados principales del capítulo]
\begin{itemize}
  \item \textbf{Anillo de vórtices.} $V_2/v_h=2{,}04$ con drogue y $1{,}98$ sin él. Ese
        valor está a \SI{1.8}{m/s} de un borde de la región a $1{,}5\,v_h$, pero coincide
        con el borde más conservador ($\approx2\,v_h$). La transición desde
        \SI{13.4}{m/s} recorre el plano $V/v_h$ de 4,5 a 2,0 sin entrar en la región.
  \item \textbf{Péndulo.} Hay un modo de inclinación del cuerpo de
        \SIrange{2.3}{2.6}{Hz} con $\zeta\approx0{,}22$, partido en dos ramas por la
        giroscopía. El modo del sistema drogue--cuerpo es de \SI{0.18}{Hz} con
        $\zeta\approx0{,}84$. Ningún autovalor cruza al semiplano derecho en los barridos
        de centro de masa, rpm, largo de cuerda y fuerza H.
  \item \textbf{Giro del cuerpo.} El par de los rodamientos es de
        $\approx\SI{0.73}{mN.m}$ (rango \numrange{0.34}{1.7}). Sin aletas el cuerpo llega a
        $\approx\num{160}$~rpm a los \SI{10}{s} y a $\approx\num{540}$~rpm a los
        \SI{40}{s}. Cuatro aletas desplegables de $40\times\SI{40}{mm}$ lo dejan en
        $\approx\SI{2.5}{rad/s}$. Para cumplir el criterio de video
        ($\le\SI{0.5}{rad/s}$) hay que darles un calaje de $\approx\ang{1.3}$ contra el
        par y ajustarlo en ensayo.
  \item \textbf{Batimiento.} $\nu_\beta=1{,}20$ y $\zeta_\beta=0{,}60$: el modo está
        muy amortiguado. Un $\delta_3$ de \ang{15} a \ang{20} baja la conicidad a
        $\approx\ang{4.5}$ y aleja la pala del tope, pero agrega $\approx\ang{-1.5}$ de
        paso.
  \item \textbf{Divergencia y flameo.} La sección típica da para la pala A
        $\Omega_D/\Omega=0{,}96$ (cota conservadora). El modelo p-k del
        capítulo~\ref{cap:estructura}, reescalado, da un margen de flameo de 1,53 y uno de
        divergencia de 1,19 a 1,59. La pala BV queda con márgenes de 3,2 a 4,2.
        \textbf{La pala BV es necesaria.}
  \item \textbf{Drogue.} Para tocar las puntas, la cuerda tendría que inclinarse
        \ang{86} respecto del eje. Eso solo pasa si el drogue pierde tensión en la estela
        del rotor. Se recomienda un tubo guía rígido sobre el cubo.
\end{itemize}
\end{resultado}

%=====================================================================
\section{Anillo de vórtices}\label{sec:estabilidad-vrs}

El punto de diseño está fuera de la región de anillo de vórtices, pero en su borde
superior. La velocidad inducida en vuelo estacionario con el empuje del punto de diseño
es
\begin{equation}
  v_h=\sqrt{\frac{T}{2\rho A}}=\sqrt{\frac{3{,}58}{2\cdot1{,}225\cdot0{,}1257}}
  =\SI{3.41}{m/s},
  \qquad \frac{V_2}{v_h}=2{,}04 .
  \label{eq:estabilidad-vh}
\end{equation}
Sin drogue, $T=W$, $v_h=\SI{3.99}{m/s}$ y $V_2/v_h=1{,}98$. El cociente es casi el mismo
porque, con menos empuje, la velocidad de descenso y $v_h$ bajan juntas.

Los límites publicados de la región no coinciden entre sí. Dependen del criterio:
fluctuación de empuje, pérdida de control o cambio del flujo medio. Para descenso axial,
la región de flujo muy inestable se sitúa aproximadamente entre $0{,}4$ y $1{,}5\,v_h$.
Su borde superior se lleva a $\approx2\,v_h$ con criterios más amplios, y por encima está
el estado de molino turbulento~\cite{leishman2006,johnson2005}. Las fronteras de la
figura~\ref{fig:estabilidad-vrs} son esos valores redondeados, no curvas medidas de este
rotor. Toda autorrotación axial real trabaja cerca de $V/v_h\approx1{,}7$--$2$: el rotor
funciona como un molino turbulento y es inevitable estar cerca del borde.

\begin{figure}[H]
\centering
\begin{tikzpicture}
\begin{axis}[width=0.9\textwidth, height=6.2cm, xmin=0, xmax=6, ymin=0, ymax=5,
  xlabel={Tiempo desde la liberación [s]}, ylabel={$V/v_h$}, grid=major,
  legend pos=north east, legend style={font=\footnotesize}]
  \fill[cfreno!25] (axis cs:0,0.4) rectangle (axis cs:6,1.5);
  \fill[cfreno!10] (axis cs:0,1.5) rectangle (axis cs:6,2.0);
  \addplot[crot, very thick] table[x=t, y=r] {analysis/data/estabilidad_vrs.dat};
  \addlegendentry{trayectoria (BEMT, transitorio)}
  \addplot[cfijo, dashed, thick] coordinates {(0,1.98) (6,1.98)};
  \addlegendentry{sin drogue, 1,98}
  \node[font=\scriptsize, cfreno!70!black] at (axis cs:3,0.95) {anillo de vórtices ($0{,}4$--$1{,}5\,v_h$)};
  \node[font=\scriptsize, cfreno!70!black] at (axis cs:1.6,1.75) {borde amplio ($\le 2\,v_h$)};
  \node[font=\scriptsize, anchor=west] at (axis cs:4.2,2.45) {diseño: 2,04};
\end{axis}
\end{tikzpicture}
\caption{Transición de \SI{13.4}{m/s} al equilibrio en el plano $V/v_h$. Se usa el
transitorio del capítulo~\ref{cap:bemt}, con $v_h(t)=\sqrt{T(t)/2\rho A}$. Las bandas son
límites aproximados de la literatura.}
\label{fig:estabilidad-vrs}
\end{figure}

\paragraph{Transición.} Durante el arranque el empuje es bajo ($T=\SI{2.7}{N}$ con
$\Omega=0$) y la velocidad es alta. $V/v_h$ arranca en 4,5 y baja de forma monótona hasta
2,04 en unos \SI{6}{s}, sin sobrepasarse. El rotor llega a la zona de anillo por arriba y
nunca la cruza.

\paragraph{Margen ante ráfagas y caída de rpm.} Una ráfaga ascendente $w$ reduce la
velocidad relativa a $V_2-w$. Para llegar a $1{,}5\,v_h$ hace falta
$w=V_2-1{,}5v_h=\SI{1.8}{m/s}$. Para el borde de $2\,v_h$ basta con \SI{0.12}{m/s}, de
modo que contra ese criterio el margen es nulo. La perturbación dura poco: la constante
de tiempo vertical es $\tau_V=m/(\partial D/\partial V)\approx mV_2/2W=\SI{0.35}{s}$. Una
caída de rpm baja el empuje, el conjunto acelera y $V/v_h$ sube: el sistema se aleja
solo de la región. La caída de rpm, por lo tanto, no lleva al anillo de vórtices.
Además, $\partial T/\partial V>0$ en el molino turbulento, lo que amortigua el
movimiento vertical.

\begin{nota}[Interpretación]
Estar a $\approx2\,v_h$ significa trabajar en molino turbulento, con fluctuaciones de
empuje esperables de algunos por ciento, no dentro del anillo de vórtices. El dato
crítico es el borde a $1{,}5\,v_h$, con un margen de \SI{1.8}{m/s} (\SI{26}{\percent}).
Un $C_R$ mayor (Buhl, $1{,}25$) llevaría a $V_2/v_h\approx1{,}8$ y reduciría ese margen a
$\approx\SI{1}{m/s}$.
\end{nota}

%=====================================================================
\section{Dinámica de péndulo y giroscópica}\label{sec:estabilidad-pend}

\subsection{Momento angular y su transmisión al cuerpo}
Con bisagras de batimiento, el cuerpo apenas siente el momento angular del rotor. La
inercia polar del rotor es
$I_p=B(I_b+2eS_b+m_be^2)\approx\SI{3.9e-4}{kg.m^2}$, con $S_b=m_bx_g=\SI{4.37e-4}{kg.m}$.
Su momento angular es $H=I_p\Omega=\SI{0.061}{N.m.s}$. Si el rotor fuera rígido, una
inclinación del eje a velocidad $q$ exigiría un par $Hq$. Con bisagras, el disco se
inclina con retraso y solo el resorte efectivo de la excentricidad,
\begin{equation}
  K_h=\frac{B}{2}eS_b\Omega^2=\SI{0.95}{N.m/rad},
  \label{eq:estabilidad-Kh}
\end{equation}
transmite par al eje. Con la respuesta cuasiestacionaria del plano de puntas,
$\beta_\text{rel}=-(16/\gamma+i)\,\dot\theta/\Omega$~\cite{johnson1980}, el cuerpo recibe
un amortiguamiento $K_h\,16/(\gamma\Omega)=\SI{8.4e-3}{N.m.s}$ y un término
giroscópico $H_\text{ef}=H_\text{cubo}+K_h/\Omega=\SI{6.1e-3}{N.m.s}$. Es el
\SI{10}{\percent} de $H$, de modo que la precesión del conjunto es débil y no domina.

\subsection{Modelo linealizado}
Se usa notación compleja ($X=x+iy$, $\theta=\theta_x+i\theta_y$) porque el conjunto es
axisimétrico. Los grados de libertad son el desplazamiento lateral del centro de masa
$X$, la inclinación $\theta$ y el desplazamiento lateral del drogue $Y$:
\begin{align}
  m\ddot X &= -c_r(\dot X+h\dot\theta)+(1-k_H)T\theta+\frac{N}{L}(Y-X-h_a\theta),\\
  I_t\ddot\theta &= iH_\text{ef}\dot\theta-\Big(c_rh^2+\frac{16K_h}{\gamma\Omega}\Big)\dot\theta
   -c_rh\dot X-k_HTh\,\theta+h_a\frac{N}{L}(Y-X-h_a\theta)-Nh_a\theta,\\
  m_d\ddot Y &= -\frac{N}{L}(Y-X-h_a\theta)-\frac{2N}{V}\dot Y .
  \label{eq:estabilidad-modelo}
\end{align}
Aquí $c_r=k_HT/V$ es la derivada de fuerza H del rotor: con $k_H=1$ la resultante del
rotor se alinea con el viento relativo, como en un paracaídas. El factor $k_H$ se barre
entre 0,3 y 1,5 porque no se conoce bien. El empuje va a lo largo del eje y pasa por el
centro de masa, así que no produce momento. La rigidez de péndulo viene de la fuerza H en
el cubo y de la tensión $N$ de la cuerda en su salida, por encima del centro de masa.
Valores nominales: $h=\SI{0.15}{m}$ (cubo sobre el centro de masa, est.),
$h_a=h+\SI{0.03}{m}$, $I_t=\SI{3e-3}{kg.m^2}$ (est.), $L=\SI{1.5}{m}$,
$N=W-T-D_\text{cuerpo}=\SI{1.16}{N}$ y $m_d=\SI{15}{g}$ (drogue más masa aparente,
est.). El cuerpo cilíndrico no aporta fuerza lateral lineal: su fuerza de flujo cruzado
es cuadrática.

\begin{table}[H]
\centering\small
\caption{Autovalores del modelo nominal. El modo de traslación lateral pura ($\lambda=0$)
es neutro, porque nada fija la posición horizontal.}
\label{tab:estabilidad-modos}
\begin{tabular}{@{}lcccc@{}}
\toprule
Modo & $\lambda$ [1/s] & $f$ [Hz] & $\zeta$ & Sin drogue\\\midrule
Inclinación, rama progresiva & $-3{,}60\pm16{,}6i$ & 2,64 & 0,21 & $-3{,}97\pm16{,}1i$\\
Inclinación, rama regresiva  & $-3{,}18\pm14{,}6i$ & 2,32 & 0,21 & $-3{,}46\pm14{,}0i$\\
Péndulo drogue--cuerpo       & $-1{,}74\pm1{,}14i$ & 0,18 & 0,84 & ---\\
Lateral del drogue           & $-19{,}7$ & --- & sobreamort. & ---\\
Deriva lateral amortiguada   & --- & --- & --- & $-1{,}27$\\
\bottomrule
\end{tabular}
\end{table}

\subsection{Lugar de raíces}
Ningún barrido lleva un autovalor al semiplano derecho. Al bajar el centro de masa (más
$h$) sube la frecuencia del modo de inclinación, porque crece la rigidez $k_HTh$, y
también su amortiguamiento, que crece con $c_rh^2$. Con $h=\SI{0.05}{m}$ el modo baja a \SIrange{1.35}{1.7}{Hz} y $\zeta=0{,}17$.  Las rpm actúan a través de $T\propto\Omega^2$ y de $K_h$.
El largo de cuerda casi no cambia el modo de inclinación: solo mueve el péndulo lento.

\begin{figure}[H]
\centering
\begin{tikzpicture}
\begin{groupplot}[group style={group size=2 by 2, horizontal sep=1.6cm, vertical sep=1.7cm},
  width=0.47\textwidth, height=5.2cm, grid=major, xlabel={$\mathrm{Re}\,\lambda$ [1/s]},
  ylabel={$\mathrm{Im}\,\lambda$ [rad/s]}, xmin=-8, xmax=0.5, ymin=-20, ymax=20,
  colorbar style={font=\scriptsize}, point meta=explicit]
\nextgroupplot[title={\footnotesize Centro de masa $h$ = 0,05 a \SI{0.22}{m}}]
  \addplot[only marks, mark=*, mark size=1.1pt, scatter, colormap/viridis] table[x=re, y=im, meta=p] {analysis/data/estabilidad_rl_h.dat};
\nextgroupplot[title={\footnotesize $\Omega/\Omega_0$ = 0,5 a 1,2}]
  \addplot[only marks, mark=*, mark size=1.1pt, scatter, colormap/viridis] table[x=re, y=im, meta=p] {analysis/data/estabilidad_rl_om.dat};
\nextgroupplot[title={\footnotesize Largo de cuerda $L$ = 0,5 a \SI{3}{m}}]
  \addplot[only marks, mark=*, mark size=1.1pt, scatter, colormap/viridis] table[x=re, y=im, meta=p] {analysis/data/estabilidad_rl_L.dat};
\nextgroupplot[title={\footnotesize Factor de fuerza H $k_H$ = 0,3 a 1,5}]
  \addplot[only marks, mark=*, mark size=1.1pt, scatter, colormap/viridis] table[x=re, y=im, meta=p] {analysis/data/estabilidad_rl_kH.dat};
\end{groupplot}
\end{tikzpicture}
\caption{Lugar de raíces del modelo~\eqref{eq:estabilidad-modelo}. El color va de oscuro
(valor bajo del parámetro) a claro (valor alto). El modo lateral del drogue
($\approx-20$) queda fuera de escala.}
\label{fig:estabilidad-rl}
\end{figure}

\begin{nota}[Límites del modelo]
El modelo es cuasiestacionario: el disco sigue al eje sin dinámica propia y la estela no
tiene retraso. La fuerza H ($k_H$), $I_t$ y $h$ son estimaciones. Como criterio se exige
$\mathrm{Re}\,\lambda<0$ con $\zeta\ge0{,}1$ en todos los modos oscilatorios dentro del
rango barrido. Se cumple con margen. La amplitud de balanceo ante turbulencia no se
calcula: la oscilación de \SIrange{2.3}{2.6}{Hz} es la que verá la cámara nadir.
\end{nota}

%=====================================================================
\section{Giro del cuerpo y aletas antigiro}\label{sec:estabilidad-giro}

Sin aletas, el cuerpo gira cada vez más rápido durante todo el descenso. La causa es el
par de rozamiento de los rodamientos, que arrastra al cuerpo en el sentido del rotor.

\paragraph{Par de rozamiento.} Por rodamiento se suman un término de carga y uno viscoso.
El de carga sigue el modelo simplificado de SKF, $M_1=\tfrac12\mu Pd$, con
$\mu=0{,}0015$, $P=\SI{12.1}{N}$ (tabla~\ref{tab:estructura-rod}) y $d=\SI{8}{mm}$. El
viscoso sigue a Palmgren, $M_0=10^{-7}f_0(\nu n)^{2/3}d_m^3$ [N\,mm], con
$d_m=\SI{12}{mm}$ y $n$ en rpm. Para grasa se toman valores típicos (no de catálogo):
$f_0=1{,}5$ (de 0,7 a 2) y $\nu=\SI{25}{mm^2/s}$ (de 15 a 70). El escudo 2Z no roza.
Resulta
\begin{equation}
  M_f=2(M_0+M_1)\approx\SI{0.73}{mN.m}\qquad(\text{rango }0{,}34\text{--}\SI{1.7}{mN.m}),
\end{equation}
dominado por la grasa. Depende de la velocidad relativa $\Omega-\omega_b$.

\paragraph{Amortiguamiento del cuerpo.} El roce de la piel del cilindro (capa laminar,
$C_f=1{,}33/\sqrt{Re_L}$) da $c_b=\SI{2.8e-6}{N.m.s}$, casi nada. Con
$I_z=\SI{4e-4}{kg.m^2}$ (est.), la aceleración inicial es $M_f/I_z\approx\SI{1.8}{rad/s^2}$.
El cuerpo llega a \num{160}~rpm a los \SI{10}{s} y a \num{540}~rpm a los \SI{40}{s}
(figura~\ref{fig:estabilidad-giro}).

\paragraph{Criterio para el video.} Se pide $|\omega_b|\le\SI{0.5}{rad/s}$
($\approx\SI{29}{\degree/s}$). Así la imagen gira menos de \ang{1} por cuadro a
\SI{30}{fps} y se puede estabilizar después. El criterio es propio, no del concurso.

\paragraph{Aletas.} El envolvente Ø\SI{95}{mm} solo deja \SI{3.5}{mm} radiales. Unas
aletas fijas de esa altura darían $c\sim\num{e-6}$, inútiles. Las aletas tienen que ser
\textbf{desplegables}: plegadas contra el cuerpo y abiertas por resorte al liberar el
rotor. Cada aleta de envergadura $s$ y cuerda $l$, con centro a $r_f=R_c+s/2$, aporta
\begin{equation}
  c_f=\tfrac12\rho V a_fSr_f^2,\qquad
  M_\delta=\tfrac12\rho V^2a_fSr_f\,\delta,\qquad
  a_f=\frac{2\pi\mathcal A}{2+\sqrt{\mathcal A^2+4}},\quad\mathcal A=2s/l .
\end{equation}
Con cuatro aletas de $40\times\SI{40}{mm}$ ($a_f=2{,}6$): $c=\SI{2.9e-4}{N.m.s}$ y el
cuerpo se estabiliza en $M_f/c\approx\SI{2.5}{rad/s}$ (24~rpm). Solo con amortiguamiento
haría falta $s\approx\SI{80}{mm}$. Es más eficaz un \textbf{calaje} $\delta$ que anule el
par: $\delta=M_f/M_\delta\approx\ang{1.3}$. El resto se debe a la incertidumbre de $M_f$:
con un error de $\pm\SI{50}{\percent}$ quedan $\approx\SI{1.2}{rad/s}$. Por eso el calaje
se ajusta midiendo el giro en un ensayo de túnel o de caída.

\begin{figure}[H]
\centering
\begin{tikzpicture}
\begin{groupplot}[group style={group size=2 by 1, horizontal sep=1.8cm},
  width=0.47\textwidth, height=5.4cm, grid=major]
\nextgroupplot[xlabel={Tiempo [s]}, ylabel={Giro del cuerpo [rpm]}, ymode=log, ymin=1, ymax=1000,
  legend pos=south east, legend style={font=\scriptsize}]
  \addplot[cfreno, very thick] table[x=t, y=w0] {analysis/data/estabilidad_giro.dat};
  \addlegendentry{sin aletas}
  \addplot[cfijo, very thick] table[x=t, y=w1] {analysis/data/estabilidad_giro.dat};
  \addlegendentry{4 aletas 40×40}
  \addplot[ccuerda, very thick, dashed] table[x=t, y=w2] {analysis/data/estabilidad_giro.dat};
  \addlegendentry{4 aletas 60×50}
\nextgroupplot[xlabel={Envergadura de aleta $s$ [mm]}, ylabel={$\omega_b$ [rad/s] / calaje [\si{\degree}]},
  ymode=log, ymin=0.3, ymax=60, legend style={font=\scriptsize}, legend pos=north east]
  \addplot[cfijo, very thick] table[x=s, y=w1] {analysis/data/estabilidad_aletas.dat};
  \addlegendentry{$\omega_b$ sin calaje}
  \addplot[crot, very thick, dashed] table[x=s, y=cant] {analysis/data/estabilidad_aletas.dat};
  \addlegendentry{calaje para $\omega_b=0$}
  \addplot[cgris, thick, densely dotted] coordinates {(10,0.5) (80,0.5)};
\end{groupplot}
\end{tikzpicture}
\caption{(a) Giro del cuerpo durante la etapa 2 con $M_f$ nominal. (b) Cuatro aletas de
cuerda \SI{40}{mm}: giro de equilibrio sin calaje y calaje que anula el par. La línea de
puntos es el criterio de \SI{0.5}{rad/s}.}
\label{fig:estabilidad-giro}
\end{figure}

Las aletas también suman amortiguamiento lateral bajo el centro de masa: el modo de
péndulo lento pasa de $\zeta=0{,}84$ a 0,83 y su $\mathrm{Re}\,\lambda$ de $-1{,}74$ a
$-1{,}80$. Es un efecto chico pero favorable.

%=====================================================================
\section{Batimiento y acoplamiento paso--batimiento}\label{sec:estabilidad-flap}

El modo de batimiento está muy amortiguado. La frecuencia propia en el sistema giratorio,
con la excentricidad y el resorte de apertura, es
\begin{equation}
  \nu_\beta^2=1+\frac{eS_b}{I_b}+\frac{k_\beta}{I_b\Omega^2}=1{,}447,\qquad
  \nu_\beta=1{,}20,\qquad \zeta_\beta=\frac{\gamma}{16\nu_\beta}=0{,}60,
  \label{eq:estabilidad-nub}
\end{equation}
con autovalores $s/\Omega=-0{,}72\pm0{,}96i$. Una perturbación se reduce a $e^{-2\pi\cdot0{,}72}$ (\SI{1}{\percent})
en una vuelta. En el sistema fijo, los modos regresivo y progresivo quedan a
$|0{,}96\mp1|\,\Omega$, es decir $0{,}04\Omega$ y $1{,}96\Omega$. El regresivo es lento, pero
tiene el mismo amortiguamiento.

\paragraph{$\delta_3$.} Si el eje de batimiento se inclina $\delta_3$, el paso cambia en
$\Delta\theta=-\tan\delta_3\,\beta$. Al primer orden, eso equivale a una rigidez
aerodinámica adicional, $\nu_e^2=\nu_\beta^2+(\gamma/8)\tan\delta_3$. La conicidad baja
como $\beta_0\propto1/\nu_e^2$ (tabla~\ref{tab:estabilidad-d3}). Con $\delta_3=\ang{20}$, $\beta_0$ pasa de
\ang{6.1} a \ang{4.5}. Con la misma ráfaga del capítulo~\ref{cap:estructura}, la pala
llega a $\approx\ang{6.9}$ y queda lejos del tope de \ang{+15}. A cambio, el paso efectivo
baja \ang{1.6} y las rpm suben, por la tendencia de la figura~\ref{fig:bemt-barrido-c}.
Un $\delta_3$ negativo es peligroso: con \ang{-30} la conicidad llega a \ang{14} y la pala
apoya en el tope. Con el tope a \ang{+15} y $\delta_3=0$ el margen es suficiente, así que
el $\delta_3$ queda como opción y no como requisito.

\begin{table}[H]
\centering\small
\caption{Efecto de $\delta_3$ (estimación de primer orden).}
\label{tab:estabilidad-d3}
\begin{tabular}{@{}ccccc@{}}
\toprule
$\delta_3$ & $\nu_e$ & $\zeta$ & $\beta_0$ & $\Delta\theta$ en $\beta_0$\\\midrule
\ang{-30} & 0,79 & 0,92 & \ang{14.3} & $+\ang{8.3}$\\
\ang{-10} & 1,09 & 0,66 & \ang{7.4} & $+\ang{1.3}$\\
\ang{0}   & 1,20 & 0,60 & \ang{6.1} & 0\\
\ang{10}  & 1,30 & 0,55 & \ang{5.2} & $-\ang{0.9}$\\
\ang{20}  & 1,40 & 0,51 & \ang{4.5} & $-\ang{1.6}$\\
\ang{30}  & 1,51 & 0,48 & \ang{3.9} & $-\ang{2.2}$\\
\bottomrule
\end{tabular}
\end{table}

%=====================================================================
\section{Divergencia y flameo de la pala}\label{sec:estabilidad-flameo}

A \num{1505}~rpm, la pala A no tiene margen aeroelástico; la pala BV sí. En la sección
típica a $0{,}75R$, con el centro de corte a $x_{ac}$ detrás del centro aerodinámico y una
frecuencia de torsión
$\omega_\theta^2=\omega_{\theta0}^2+\Omega^2$ (con
$\omega_{\theta0}=\tfrac{\pi}{2L}\sqrt{GJ/I_\theta}$ y el término de hélice), la divergencia
ocurre cuando
\begin{equation}
  \Omega_D^2=\frac{I_\theta\,\omega_{\theta0}^2}{\tfrac12\rho r^2ca\,x_{ac}-I_\theta}.
  \label{eq:estabilidad-div}
\end{equation}
Con las propiedades de la tabla~\ref{tab:estructura-aeroel}, la pala A da
$\Omega_D=\SI{151}{rad/s}$, o sea $\Omega_D/\Omega=0{,}96$, y la pala BV da
$\Omega_D/\Omega=3{,}2$. La sección típica es conservadora: carga toda la pala con la
presión dinámica de $0{,}75R$ contra la rigidez de la raíz. El análisis distribuido del
capítulo~\ref{cap:estructura} da 1790--\num{2400}~rpm para A, que a las rpm refinadas son
$\Omega_D/\Omega=1{,}19$--$1{,}59$.

Para el flameo, el modelo cuasiestacionario de dos grados de libertad no sirve: sin la
función de Theodorsen predice inestabilidades débiles de torsión a cualquier velocidad.
Se adoptan entonces los resultados p-k del capítulo~\ref{cap:estructura}, que dependen
de $\Omega$ y no de la relación de diseño. Para A, \num{2310}~rpm dan
$\Omega_F/\Omega=1{,}53$. Para BV, \num{6250}~rpm dan 4,15. El margen usual en rotores es
$\ge1{,}3$ a la velocidad máxima (est.). A sobreviviría en régimen nominal, pero un
embalamiento al \SI{120}{\percent}, por ejemplo con paso de \ang{-6} y \num{1840}~rpm,
la deja a $\Omega_F/\Omega=1{,}26$ y con divergencia posible.

\begin{resultado}[Aeroelasticidad]
Hay que usar la pala BV (varilla de borde de ataque y capas a $\pm\ang{45}$). Su centro
de masa debe quedar en $y_{cm}\le0{,}41c$. Con BV los márgenes son
$\Omega_D/\Omega\ge3{,}2$ y $\Omega_F/\Omega\approx4{,}2$.
\end{resultado}

%=====================================================================
\section{Drogue en la estela del rotor}\label{sec:estabilidad-drogue}

El drogue solo puede tocar las palas si pierde la tensión. La cuerda sale del eje hueco a
unos \SI{30}{mm} sobre el plano de bisagras (est.). Las puntas, con conicidad de
\ang{6.1}, suben $(R-e)\sin\beta_0=\SI{17}{mm}$. Para que la cuerda llegue al disco
tendría que formar
\begin{equation}
  \phi_\text{crit}=90^\circ-\arctan\frac{z_\text{sal}-z_\text{punta}}{R}\approx\ang{86}
\end{equation}
con el eje, es decir, el drogue tendría que estar casi a la altura del rotor. Con tensión
$N=\SI{1.16}{N}$, una ráfaga lateral $w$ inclina la cuerda unos $\arctan(w/V_2)$:
\ang{16} con \SI{2}{m/s}. El modo de péndulo del drogue tiene $\zeta=0{,}84$, sin
amplificación. Los ángulos normales están lejos de \ang{86}.

\begin{figure}[H]
\centering
\begin{tikzpicture}[scale=1.4]
  \draw[fijo, fill=ffijo] (-0.33,-1.6) rectangle (0.33,0);
  \draw[rot] (-0.15,0) rectangle (0.15,0.14);
  \draw[rot, very thick] (0.15,0.07) -- (2.0,0.27);
  \draw[rot, very thick] (-0.15,0.07) -- (-2.0,0.27);
  \draw[eje] (0,-1.8) -- (0,2.3);
  \draw[cuerda, very thick] (0,0.3) -- (0.62,2.25) node[above, font=\scriptsize] {drogue};
  \draw[cuerda, dashed] (0,0.3) -- (2.2,0.45);
  \draw[cota] (0,0.07) -- (0,0.3);
  \node[font=\scriptsize, anchor=east] at (-0.05,0.42) {salida};
  \draw (0,1.0) arc (90:72:0.7);
  \node[font=\scriptsize] at (0.25,1.15) {\ang{16}};
  \node[font=\scriptsize, anchor=west] at (2.2,0.45) {$\phi_\text{crit}\approx\ang{86}$};
  \node[font=\scriptsize, anchor=west] at (2.0,0.15) {punta, $\beta_0=\ang{6.1}$};
  \draw[cfreno, thick, -{Stealth}] (1.8,1.5) -- (1.8,2.1) node[right, font=\scriptsize, black] {aire relativo};
\end{tikzpicture}
\caption{Geometría de interferencia (sin escala). La cuerda tendría que bajar a la línea
discontinua para alcanzar las puntas.}
\label{fig:estabilidad-drogue}
\end{figure}

El riesgo real está en la estela. El aire relativo sube y el drogue queda aguas abajo del
rotor, en una estela de molino turbulento con poco flujo medio y mucha recirculación. A
$L=7{,}5R$ la estela ya se mezcló en parte, pero la presión dinámica sobre el drogue
fluctúa y puede caer. Si el drogue pierde tensión, el péndulo pierde rigidez, porque
$N\to0$ en la ec.~\eqref{eq:estabilidad-modelo}, y la tela puede caer hacia el cubo. Las
prestaciones no dependen de esto, porque el caso sin drogue ($V_2=\SI{7.91}{m/s}$)
cumple C4. Durante la transición el cuerpo frena a $\approx\SI{3}{m/s^2}$. El drogue
puede frenar a $N/m_d\approx\SI{80}{m/s^2}$, así que la cuerda sigue tensa. No se
modela la estela en el drogue, y este riesgo queda como \textbf{cualitativo}.

%=====================================================================
\section{Márgenes y recomendaciones}\label{sec:estabilidad-rec}

\begin{table}[H]
\centering\small
\caption{Márgenes de estabilidad en el punto de diseño.}
\label{tab:estabilidad-margenes}
\begin{tabularx}{\textwidth}{@{}lYYc@{}}
\toprule
Fenómeno & Valor & Criterio & Estado\\\midrule
Anillo de vórtices & $V/v_h=2{,}04$; ráfaga \SI{1.8}{m/s} & $>1{,}5\,v_h$ & justo\\
Inclinación (péndulo) & $\zeta=0{,}21$, \SI{2.5}{Hz} & $\zeta\ge0{,}1$ & cumple\\
Péndulo drogue & $\zeta=0{,}84$ & $\zeta\ge0{,}1$ & cumple\\
Batimiento & $\zeta_\beta=0{,}60$, $\nu_\beta=1{,}20$ & amortiguado & cumple\\
Conicidad + ráfaga & \ang{9.3} & $<$ tope \ang{15} & cumple\\
Giro del cuerpo, sin aletas & 160--540 rpm & $\le\SI{0.5}{rad/s}$ & no cumple\\
Giro, 4 aletas 40×40 + calaje & $\approx0$--\SI{1.2}{rad/s} & $\le\SI{0.5}{rad/s}$ & con ajuste\\
Divergencia, pala A / BV & 0,96--1,59 / $\ge3{,}2$ & $\ge1{,}3$ & A no, BV sí\\
Flameo, pala A / BV & 1,53 / 4,15 & $\ge1{,}3$ & BV\\
Drogue contra palas & \ang{16} frente a \ang{86} & con tensión & cualitativo\\
\bottomrule
\end{tabularx}
\end{table}

\begin{enumerate}
  \item \textbf{Pala BV}, con el centro de masa en $\le0{,}41c$ y verificado con el
        ensayo de balanceo sobre filo.
  \item \textbf{Cuatro aletas desplegables} de $\approx40\times\SI{40}{mm}$ en el fondo
        del cuerpo, abiertas por resorte al liberar el rotor, con calaje inicial de
        \ang{1.3} contra el par del rotor y ajustable entre \ang{0} y \ang{3}. Medir el
        par de los rodamientos en banco (caída libre de un volante) antes de fijar el
        calaje. Reengrasar con aceite ligero o desengrasar los 688ZZ reduce $M_f$ a
        menos de la mitad.
  \item \textbf{Centro de masa bajo}: lastre en el fondo, con
        $h\gtrsim\SI{0.12}{m}$ bajo el cubo, para mantener la frecuencia y el
        amortiguamiento del modo de inclinación.
  \item \textbf{Tubo guía} rígido de \SIrange{40}{60}{mm} sobre el cubo para la cuerda, y
        destorcedor por encima. Mantener $L\ge\SI{1.5}{m}$.
  \item $\delta_3=\ang{15}$--\ang{20} en las bisagras solo si los ensayos muestran
        conicidades mayores que \ang{7}. Si se usa, corregir el paso para mantener las
        rpm.
  \item Verificar en ensayo de caída: rpm, $V_2$, balanceo (IMU) y giro del cuerpo. Ningún
        valor de este capítulo está medido.
\end{enumerate}
